\section{Introduction}

With the rise of tools like ChatGPT, HuggingFace and MidJourney, the need for efficient image generation tools has drastically increased over the last 5 years. The traditionally superior Generative Adversarial Networks (GAN) \cite{goodfellow2014generativeadversarialnetworks, DBLP:journals/corr/abs-1809-11096, DBLP:journals/corr/abs-1912-00953} and Variational Auto-Encoder (VAE) suffer from a major drawback : they are hard to train and require careful hyper-parameter selection to avoid collapse during training \cite{DBLP:journals/corr/abs-1802-05957}. In comparison, Diffusion Models  \cite{DBLP:journals/corr/abs-2006-11239} are more stable and scalable, but they require hundreds of steps to generate a clean image: training time is traded against sampling time. Moreover, in large image datasets such as IMAGENET-64, they are not competitive with GANs \cite{DBLP:journals/corr/abs-2102-09672}. More recently, Consistency Models \cite{song2023consistencymodels, yang2024consistencyflowmatchingdefining, geng2024consistencymodelseasy} and Flow Map Matching \cite{boffi2025flowmapmatchingstochastic, peng2025flowanchoredconsistencymodels} methods aim at generating images in only 1 or 2 steps and significantly improve sampling performances. \\

Diffusion Models generate images by applying successive denoising steps to the input signal, until its distribution matches this of the original data. In practice, it reverses the noising process, which is solving a Differential Equations (DE). Other approaches that aim at solving DE using data-driven approaches employ Neural Operators \cite{ganeshram2025fcpinohighprecisionphysicsinformed, DBLP:journals/corr/abs-2108-08481} : Neural Network composed of several Spectral Convolution layers is trained to approximate the solution to the DE by projecting the input in the frequency domain, heavily relying on the Fast Fourier Transform algorithm. Interestingly, they are rarely considered for improving diffusion models; we only found \cite{zheng2023fastsamplingdiffusionmodels} that propose a UNET architecture composed of Fourier Neural Operator blocks to generate images. \\

In this work, we propose to enhance Consistency Models by employing a Neural Operator architecture. 
In Section~\ref{sec:diffno}, we detail the theoretical framework that we use for Consistency Models and Neural Operators.  
In Section~\ref{sec:stno}, we describe our approach to train a NO for image diffusion. We discuss the choice of the distance metric and we propose a Self-Training Loss that allows the model to train without the need of a pre-trained solver.
In Section~\ref{}, we give our experimental set-up, and evaluate our model on CIFAR-10 and IMAGENET-64. We perform an ablation study to analyze the critical parts of our approach. Our code is available on github : \repo. 
In Section~\ref{}, we discuss the future directions of research that we aim to pursue to further improve our model.